\documentclass[11pt]{article}
\usepackage{fontspec}
\setmainfont{DejaVu Serif}
\usepackage{amsmath,amssymb}
\usepackage{geometry}\geometry{margin=2.3cm}
\usepackage{xcolor}\usepackage{graphicx}\usepackage{enumitem}
\usepackage[hidelinks]{hyperref}

\definecolor{ink}{HTML}{16212e}\definecolor{accent}{HTML}{1f5fa8}\definecolor{soft}{HTML}{eef3fb}
\color{ink}
\newcommand{\key}[1]{\par\medskip\noindent\colorbox{soft}{\parbox{\dimexpr\linewidth-2\fboxsep}{\textcolor{accent}{\textbf{Κλειδί:}}\ #1}}\par\medskip}
\newcommand{\ex}[1]{\par\smallskip\noindent\textbf{\textcolor{accent}{Παράδειγμα.}}\ \textit{#1}\par\smallskip}
\setlength{\parskip}{0.5em}\setlength{\parindent}{0pt}

\title{\textbf{Φροντιστήριο: Ένα πειθαρχημένο pipeline για αδύναμα σήματα μετοχών}\\[3pt]
\large PEAD $\cdot$ lead-lag $\cdot$ utility gate $\cdot$ risk parity $\cdot$ regime layer\\[-1pt]
\normalsize\textcolor{accent}{Από το σήμα στο χαρτοφυλάκιο — με απλά παραδείγματα}}
\author{QuantIQ / AGEL AI}\date{}

\begin{document}
\maketitle\thispagestyle{empty}

\noindent Αυτό το φυλλάδιο εξηγεί, από το μηδέν και με μικρά αριθμητικά παραδείγματα, ολόκληρο
το σύστημα που χτίσαμε: πώς βρίσκουμε αδύναμα αλλά τίμια σήματα από δημόσια δεδομένα κερδών,
πώς ελέγχουμε αν είναι αληθινά, πώς τα συνδυάζουμε, και πώς τα τυλίγουμε σε έλεγχο ρίσκου.
Δεν χρειάζεται προηγούμενη γνώση πέρα από βασική άλγεβρα.

\tableofcontents
\bigskip

%====================================================================
\section{Η μεγάλη εικόνα: γιατί είναι δύσκολο}

Στις αγορές, ό,τι φαίνεται στο διάγραμμα τιμής το βλέπουν όλοι — και έχει ήδη «τιμολογηθεί».
Άρα ένα πραγματικό \textbf{σήμα} πρέπει να φέρνει \emph{νέα πληροφορία από έξω} (κέρδη, σχέσεις
εταιρειών, κείμενο), να προβλέπει \emph{κατεύθυνση}, και να έχει \emph{μηχανισμό που αντέχει}.

Δεύτερη σκληρή αλήθεια: κανένα μεμονωμένο τέτοιο σήμα δεν είναι δυνατό. Είναι όλα \emph{αδύναμα}.
Το «μυστικό» δεν είναι ένα μαγικό σήμα — είναι ο \textbf{συνδυασμός} πολλών αδύναμων,
ασυσχέτιστων σημάτων, με πειθαρχία ώστε να μην ξεγελιόμαστε από τον θόρυβο.

\key{Το σύστημα είναι μια \emph{αλυσίδα}: σήματα $\rightarrow$ έλεγχος (gate) $\rightarrow$
συνδυασμός (risk parity) $\rightarrow$ έλεγχος ρίσκου (regime). Κανένα μοντέλο δεν «μαθαίνει
κατεύθυνση»· εφαρμόζουμε γνωστά σήματα και τα διαχειριζόμαστε σωστά.}

%====================================================================
\section{Τα μαθηματικά εργαλεία (γρήγορη επανάληψη)}

\textbf{Μέση τιμή} $\bar{x}=\frac{1}{N}\sum x_i$. \textbf{Τυπική απόκλιση}
$\sigma=\sqrt{\frac1N\sum(x_i-\bar{x})^2}$ — «πόσο κουνιέται». \textbf{z-score} $z=(x-\bar{x})/\sigma$
— κοινή κλίμακα. \textbf{Συσχέτιση} $\rho\in[-1,1]$ — πάνε μαζί; \textbf{Κλίση} $b=\sum x_iy_i/\sum x_i^2$
(όταν $\bar{x}=0$) — πόσο αλλάζει το $y$ ανά μονάδα $x$. \textbf{Εκθετική φθορά}
$y(k)=b_0e^{-k/\tau}$ με \textbf{μισή ζωή} $t_{1/2}=\tau\ln2$.

\ex{Για $2,4,9$: $\bar{x}=5$, αποκλίσεις $-3,-1,4$, $\sigma=\sqrt{26/3}\approx2{,}9$.}

%====================================================================
\section{Σήμα 1: PEAD — η ολίσθηση μετά τα κέρδη}

Όταν μια εταιρεία εκπλήσσει θετικά στα κέρδη, η μετοχή της δεν ανεβαίνει αμέσως όσο πρέπει —
\emph{συνεχίζει} να ολισθαίνει για εβδομάδες, γιατί οι επενδυτές υπο-αντιδρούν.

\textbf{Μέτρηση της έκπληξης (SUE)} — χωρίς αναλυτές, μόνο από δημόσια κέρδη:
\[ \mathrm{SUE}=\frac{E_t-E_{t-4}}{\sigma},\quad E_{t-4}=\text{ίδιο τρίμηνο πέρυσι}. \]
\ex{Φέτος $1{,}30$, πέρυσι $1{,}05$ $\Rightarrow$ έκπληξη $+0{,}25$. Αν $\sigma=0{,}055$, τότε
$\mathrm{SUE}=0{,}25/0{,}055\approx+4{,}5$ — μεγάλη θετική έκπληξη.}

Long τα υψηλά SUE, short τα χαμηλά. Είσοδος \textbf{T+1} (την επόμενη μέρα της ανακοίνωσης),
κράτημα $\sim$60 μέρες. Όλα από το \textbf{SEC EDGAR} (δωρεάν, point-in-time).

%====================================================================
\section{Πόσο κρατάει το σήμα; (event study)}

Αντί να υποθέσουμε «60 μέρες», το \emph{μετράμε}. Στοιβάζουμε χιλιάδες ανακοινώσεις στη «μέρα 0».

\textbf{Abnormal απόδοση} $=$ απόδοση μετοχής $-$ μέσος όρος αγοράς εκείνη τη μέρα (κρατάμε το
\emph{σχετικό} κομμάτι). \textbf{Προφίλ ολίσθησης} $b(k)$ = abnormal απόδοση τη μέρα $k$ μετά την
είσοδο, ανά μονάδα SUE:
\[ b(k)=\frac{\sum_e z_e\,\mathrm{AR}_{e,k}}{\sum_e z_e^2}. \]
\ex{Τρία γεγονότα $z=(+2,0,-1)$, αποδόσεις μέρα $k$: $(+0{,}008,+0{,}001,-0{,}003)$.
$b(k)=\frac{2(0{,}008)+0+(-1)(-0{,}003)}{4+0+1}=\frac{0{,}019}{5}=0{,}0038=38$ μ.β.}

\textbf{Σωρευτικό} $B(h)=\sum_{k\le h}b(k)$ ανεβαίνει και μετά επιπεδώνει (το drift εξαντλείται).
\textbf{Φθορά}: ταιριάζουμε $b(k)\approx b_0e^{-k/\tau}$ $\Rightarrow$ μισή ζωή $t_{1/2}=\tau\ln2$.

\begin{center}\includegraphics[width=\linewidth]{pead_results.png}\\
{\small Αριστερά: $b(k)$ (μπάρες) + εκθετικό fit ($t_{1/2}\!\approx\!9$ μέρες). Μέση: σωρευτικό
$B(h)$. Δεξιά: IR ανά έτος (durability). Synthetic validation: ανέκτησε half-life $8{,}6$ ενώ η
αλήθεια ήταν $10{,}4$.}\end{center}

%====================================================================
\section{Είναι αληθινό; Newey--West και IR}

Δύο παγίδες φουσκώνουν τα στατιστικά: τα 60-ήμερα παράθυρα \emph{επικαλύπτονται}, και οι
ανακοινώσεις \emph{μαζεύονται} στο earnings season. Λύση: φτιάχνουμε ένα \textbf{ημερολογιακό
χαρτοφυλάκιο} (κάθε μέρα, long/short όλα τα ενεργά γεγονότα) και κρίνουμε την ημερήσια σειρά
αποδόσεων. Η \textbf{Newey--West} διόρθωση μεγαλώνει σωστά το σφάλμα όταν οι μέρες
αυτοσυσχετίζονται — αλλιώς η απλή $t$ είναι υπερβολικά αισιόδοξη.

\textbf{Information Ratio} $\mathrm{IR}=\sqrt{252}\,\bar{r}/\sigma_r$ (απόδοση ανά ρίσκο, ετήσια).
Ρεαλιστικά «καλό»: $0{,}5$--$1{,}0$.

\textbf{Durability}: το ίδιο IR \emph{ανά έτος}· αρνητική τάση = το φαινόμενο φθίνει διαχρονικά
(το πιο σημαντικό μέγεθος — και για paper και για προϊόν).

%====================================================================
\section{Σήμα 2: lead-lag μεταξύ ομοκλαδικών εταιρειών}

Όταν η εταιρεία Α (π.χ. τσιπ) εκπλήσσει, αυτό προμηνύει καλά νέα και για την ομοκλαδική Β
(π.χ. κινητά) — αλλά η αγορά αργεί να το συνδέσει, οπότε η Β ολισθαίνει με καθυστέρηση. Σήμα:
για κάθε μετοχή, ο μέσος SUE των \emph{ομοκλαδικών} της που μόλις ανακοίνωσαν (ίδιο SIC code).

\textbf{Το πρόβλημα}: στο earnings season ανακοινώνουν όλοι μαζί, άρα το lead-lag μπερδεύεται με
το \emph{δικό} της PEAD της Β. Πρέπει να τα ξεχωρίσουμε.

\subsection{Ο σωστός διαχωρισμός: Fama--MacBeth}
\textbf{Cross-sectional} = συγκρίνω όλες τις μετοχές την \emph{ίδια} μέρα (η κίνηση της αγοράς
κόβεται). \textbf{Παλινδρόμηση} = προσαρμόζω: απόδοση $\approx a\cdot(\text{own})+b\cdot(\text{peer})$.
Το $b$ είναι η επίδραση του peer \emph{αφού κρατήσω σταθερό} το own.

\ex{Μια μέρα, 4 ομοκλαδικές:
\begin{center}\small
\begin{tabular}{lccc}\hline
 & own SUE & peer SUE & απόδοση\\\hline
Α & $+2$ & $+1$ & $+1{,}2\%$\\
\textbf{Β} & $\mathbf{0}$ & $\mathbf{+2}$ & $\mathbf{+1{,}0\%}$\\
Γ & $-1$ & $-1$ & $-0{,}8\%$\\
Δ & $+1$ & $-2$ & $-0{,}5\%$\\\hline
\end{tabular}\end{center}
Η Β έχει own $=0$ (καμία δική της έκπληξη) αλλά ανέβηκε $+1{,}0\%$ με peer $+2$. Αυτό το $+1\%$
\emph{δεν} εξηγείται από δικά της κέρδη $\Rightarrow$ καθαρή ένδειξη lead-lag. Η παλινδρόμηση το
βρίσκει αυτόματα για όλες.}

\textbf{«MacBeth»} = το επαναλαμβάνω \emph{κάθε μέρα}, βγάζω ένα $b_t$, και παίρνω τον μέσο όρο.
Αν τα $b_t$ είναι σταθερά θετικά (όχι τυχαία), το lead-lag είναι αληθινό.

\begin{center}\includegraphics[width=\linewidth]{leadlag_selftest_results.png}\\
{\small Synthetic: το «καθαρό» lead-lag (ελέγχει own) μετρήθηκε ξεχωριστά από το own-PEAD, και
ανιχνεύθηκε η διαχρονική φθορά (δεξιά). Αριστερά: σωρευτική απόδοση· μέση: own-PEAD αναφοράς.}
\end{center}

%====================================================================
\section{Ποιο σήμα αξίζει; Το utility-weighted gate}

Δεν κρίνουμε ένα σήμα από το «πόσο σωστά μαντεύει» (accuracy) αλλά από το \textbf{πόσο
χρήμα/ρίσκο ζυγίζει} — όπως σκοράρει η Jane Street στους Kaggle (utility $=$ weight $\times$
return). Πρακτικά: φτιάχνουμε το long/short χαρτοφυλάκιο του κάθε σήματος και κοιτάμε το
\emph{οικονομικό} του P\&L: IR, Newey--West $t$, και \emph{early vs late} (μήπως φθίνει).

\textbf{Κανόνας}: περνά μόνο αν IR$>0$, $|t_{\text{NW}}|\ge$ κατώφλι (π.χ. 2), και δεν καταρρέει
στο δεύτερο μισό. Έτσι κόβουμε σήματα που «φαίνονται» καλά in-sample αλλά είναι θόρυβος.

%====================================================================
\section{Πώς τα ενώνουμε: risk parity}

Πολλά αδύναμα σήματα $\rightarrow$ ένα τελικό σκορ. \textbf{Δεν} μοιράζουμε ίσο κεφάλαιο —
μοιράζουμε \textbf{ίσο ρίσκο}. Βάρος $\propto 1/\sigma$ (αντιστρόφως ανάλογο της μεταβλητότητας):
\[ w_i=\frac{1/\sigma_i}{\sum_j 1/\sigma_j}. \]
\ex{own-PEAD vol $2\%$/μέρα, peer lead-lag vol $1\%$/μέρα. $1/2=0{,}5$, $1/1=1$, άθροισμα $1{,}5$
$\Rightarrow$ \textbf{own $33\%$, peer $67\%$}. Έλεγχος: $0{,}33\times2\%=0{,}67\times1\%=0{,}67\%$
— ίσο ρίσκο.}

Τελικό σκορ $=0{,}33\,z(\text{own})+0{,}67\,z(\text{peer})$. Αν τα δύο \emph{διαφωνούν} σε μια
μετοχή, το σκορ της μετριάζεται — δύο ορθογώνιες πηγές που αλληλοελέγχονται. (Η πλήρης εκδοχή,
\emph{ERC}, λαμβάνει υπόψη και τη συσχέτιση μεταξύ σημάτων.)

%====================================================================
\section{Το φρένο ρίσκου: regime layer}

Η μεταβλητότητα, σε αντίθεση με τις αποδόσεις, είναι \emph{προβλέψιμη} (ομαδοποιείται). Άρα
ρυθμίζουμε την έκθεση ανάλογα με τον «καιρό» της αγοράς — \textbf{volatility targeting}:
\[ \text{πολλαπλασιαστής}=\frac{\text{στόχος vol}}{\text{πρόσφατη vol}}\ (\text{με όριο,
χθεσινή vol} \Rightarrow \text{χωρίς look-ahead}). \]
Σε ταραγμένα regimes μικραίνει την έκθεση· σε ήρεμα την αυξάνει. \emph{Δεν} προβλέπει
κατεύθυνση — ελέγχει ρίσκο (το μάθημα Jane Street + το DER $\theta$).

\ex{Στόχος vol $10\%$, πρόσφατη vol $25\%$ $\Rightarrow$ πολλαπλασιαστής $0{,}10/0{,}25=0{,}4$:
κρατάμε μόνο το $40\%$ της έκθεσης μέχρι να ηρεμήσει.}

%====================================================================
\section{Όλα μαζί: το engine}

\textbf{σήματα $\rightarrow$ gate $\rightarrow$ risk-parity $\rightarrow$ regime $\rightarrow$
χαρτοφυλάκιο.} Σε synthetic δοκιμή (ασθενές σήμα + κρίση), κάθε στρώμα έκανε τη δουλειά του:

\begin{center}\includegraphics[width=\linewidth]{engine_selftest_results.png}\\
{\small Αριστερά: σωρευτικό P\&L πριν/μετά το φρένο. Μέση: ο πολλαπλασιαστής έκθεσης πέφτει
στην κρίση. Δεξιά: βάρη risk parity (own $0{,}29$ / peer $0{,}71$).}\end{center}

Αποτέλεσμα δοκιμής: το gate πέρασε και τα δύο σήματα· το regime brake έκοψε το μέγιστο drawdown
από $-3{,}6\%$ σε $-0{,}55\%$ και βελτίωσε το IR ($6{,}9\!\to\!7{,}6$), μειώνοντας τη
μεταβλητότητα ($3{,}2\%\!\to\!0{,}9\%$). \emph{(Τα υψηλά νούμερα είναι artifact του synthetic —
αποδεικνύουν ότι τα εργαλεία δουλεύουν, όχι αποτέλεσμα αγοράς.)}

%====================================================================
\section{Η τίμια προσδοκία}

Στα δωρεάν daily δεδομένα, τα σήματα είναι \textbf{αδύναμα και πιθανότατα φθίνοντα}. Το πιο
πιθανό πραγματικό αποτέλεσμα είναι μέτριο ή \emph{null}. Αυτό \textbf{δεν} είναι αποτυχία: ένα
τίμιο null με σωστά Newey--West $t$-stats είναι έγκυρο, δημοσιεύσιμο αποτέλεσμα. Η αξία του
συστήματος είναι ότι (1) μετράει τίμια τι δουλεύει, (2) ελέγχει το ρίσκο, (3) συνδυάζει πολλά
αδύναμα σήματα σωστά — όχι ότι υπόσχεται μαγικό alpha.

\key{Το μικρότερο μοντέλο που κάνει τη δουλειά κερδίζει. Η χωρητικότητα (deep nets) είναι
ρίσκο, όχι αρετή — εκτός για ρίσκο/sizing και για νέες πηγές (κείμενο). Για \emph{κατεύθυνση}
σε daily τιμές, η απλότητα νικά.}

%====================================================================
\section{Σύνοψη συμβόλων}
\begin{center}\renewcommand{\arraystretch}{1.3}
\begin{tabular}{@{}ll@{}}\hline
\textbf{Σύμβολο} & \textbf{Σημασία}\\\hline
$\mathrm{SUE}$ & μέγεθος έκπληξης κερδών $=(E_t-E_{t-4})/\sigma$\\
$\mathrm{AR}$ & abnormal απόδοση (αφαιρεμένη η αγορά)\\
$b(k),B(h)$ & ολίσθηση ανά μέρα $k$ / σωρευτική ως $h$\\
$\tau,t_{1/2}$ & χρόνος φθοράς / μισή ζωή $=\tau\ln2$\\
$\mathrm{IC}$ & συσχέτιση σήματος–μελλοντικής απόδοσης\\
$\mathrm{IR}$ & απόδοση ανά ρίσκο (ετήσια)\\
$t_{\text{NW}}$ & $t$-στατιστική Newey--West (διορθωμένη)\\
$a,b$ (FM) & συντελεστές own / peer στο Fama--MacBeth\\
$w_i\propto1/\sigma_i$ & βάρη risk parity (ίσο ρίσκο)\\
πολλ/στής & regime exposure (vol targeting)\\\hline
\end{tabular}\end{center}

\bigskip
\key{Όλη η αλυσίδα με μια φράση: \emph{βρες αδύναμα σήματα από έξω (PEAD, lead-lag), έλεγξε
ποια είναι αληθινά (utility gate, Newey--West, durability), ένωσέ τα με ίσο ρίσκο (risk parity),
και τύλιξέ τα με φρένο μεταβλητότητας (regime) — τίμια, χωρίς overfitting.}}

\end{document}
